\TieudeT{PHONG TỎA VẬT LÍ 12 - VẬT LÍ NHIỆT}
\subsubsection*{PHẦN I. Thí sinh trả lời câu hỏi từ câu 1 đến câu 18. Mỗi câu hỏi thí sinh chỉ chọn một phương án}
\setcounter{ex}{0}
\begin{ex}
Các nhiệt kế thường dùng như nhiệt kế rượu, nhiệt kế thủy ngân,... được chế tạo dựa trên
\choice
{sự nở vì nhiệt của ống thủy tinh chứa chất lỏng}
{\True sự nở dài của cột chất lỏng trong ống thủy tinh}
{sự nở dài của một thanh kim loại thẳng}
{sự nở vì nhiệt của thể tích một lượng khí xác định ở áp suất không đổi}
\loigiai{}
\end{ex}
\begin{ex}
Nội năng của vật nào tăng nhiều nhất khi ta thả rơi bốn vật có cùng thể tích (đặc) từ cùng một độ cao xuống đất? Coi như toàn bộ độ giảm cơ năng chuyển hết thành nội năng của vật.
\choice
{Vật bằng cao su}
{Vật bằng nhôm}
{Vật bằng xốp}
{\True Vật bằng sắt}
\loigiai{}
\end{ex}
\begin{ex}
Mô hình động học phân tử về cấu tạo chất \textbf{không} đề cập đến nội dung nào sau đây?
\choice
{Các chất được cấu tạo từ các hạt riêng biệt là phân tử}
{Các phân tử chuyển động không ngừng}
{\True Các phân tử không tương tác với nhau}
{Nhiệt độ của vật càng cao thì tốc độ chuyển động của các phân tử cấu tạo nên vật càng lớn}
\loigiai{}
\end{ex}
\begin{ex}
Nhiệt độ không tuyệt đối $(0 K)$ là nhiệt độ mà tại đó các phân tử có
\choice
{động năng chuyển động nhiệt bằng không và thế năng tương tác giữa chúng là cực đại}
{động năng chuyển động nhiệt cực đại và thế năng tương tác giữa chúng là cực đại}
{\True động năng chuyển động nhiệt bằng không và thế năng tương tác giữa chúng là tối thiểu}
{động năng chuyển động nhiệt cực đại và thế năng tương tác giữa chúng là bằng không}
\loigiai{}
\end{ex}
\begin{ex}
\immini[thm]{Cho sơ đồ các hình thức chuyển thể như hình. Phát biểu nào sau đây là đúng?
\choice
{\textbf{(2)} là quá trình ngưng kết}
{\textbf{(1)} là quá trình nóng chảy}
{\textbf{(2)} là quá trình hóa hơi}
{\True \textbf{(3)} là quá trình ngưng tụ}}{\begin{tikzpicture}[draw=Mapcolor, very thick, scale = 0.8, line join=round, line cap=round, >=stealth]
	\draw[line width = 1] (0,0) circle(0.6) node []{Rắn};
	\draw[line width = 1] (3,0) circle(0.6) node []{Lỏng};
	\draw[line width = 1] (1.5,2.6) circle(0.6) node []{Khí};
	\draw[->,thick, line width = 1] (0.4,0.7) -- (1.1,2);
	\draw[->,thick, line width = 1] (2.2,0) -- (0.8,0);
	\draw[->,thick, line width = 1] (1.9,2) -- (2.6,0.7);
	\draw (1.5,-0.4) node []{\large \textbf{(1)}};
	\draw (0.25,1.45) node []{\large \textbf{(2)}};
	\draw (2.75,1.45) node []{\large \textbf{(3)}};
\end{tikzpicture}}
\loigiai{}
\end{ex}
\begin{ex}
Cặp nhiệt độ được chọn làm mốc trong thang đo nhiệt độ Kelvin là
\choice
{độ không tuyệt đối và nhiệt độ nước đóng băng}
{nhiệt độ nước đóng băng và nhiệt độ sôi của nước tinh khiết}
{nhiệt độ nước đóng băng và nhiệt độ điểm ba của nước}
{\True độ không tuyệt đối và nhiệt độ điểm ba của nước}
\loigiai{}
\end{ex}
\begin{ex}
Nhiệt độ cơ thể người là $37^\circ C$ sẽ tương ứng với nhiệt độ bao nhiêu trong thang đo nhiệt độ Kelvin?
\choice
{\True $310K$}
{$300K$}
{$236K$}
{$210K$}
\loigiai{}
\end{ex}
\begin{ex}
Nội năng của một hệ là
\choice
{\True tổng động năng và thế năng tương tác của các phân tử cấu tạo nên hệ}
{tổng công và nhiệt mà hệ truyền ra bên ngoài}
{tổng động năng và thế năng của hệ}
{tổng công và nhiệt mà hệ nhận được từ bên ngoài}
\loigiai{}
\end{ex}
\begin{ex}
Trong chuyển động nhiệt, các phân tử lỏng
\choice
{chuyển động tự do về mọi phía}
{\True dao động quanh vị trí cân bằng luôn luôn thay đổi}
{chuyển động hỗn loạn}
{dao động quanh vị trí cân bằng cố định}
\loigiai{}
\end{ex}
\begin{ex}
Với cùng một chất quá trình chuyển thể nào sau đây sẽ làm giảm lực tương tác giữa các phân tử nhiều nhất?
\choice
{Nóng chảy}
{Ngưng tụ}
{Đông đặc}
{\True Hóa hơi}
\loigiai{}
\end{ex}
\begin{ex}
Sự hóa hơi xảy ra trên bề mặt chất lỏng gọi là
\choice
{sự nóng chảy}
{\True sự bay hơi}
{sự đông đặc}
{sự sôi}
\loigiai{}
\end{ex}
\begin{ex}
Gọi A và Q lần lượt là tổng công và nhiệt lượng mà hệ nhận được, $\Delta U $ là độ biến thiên nội năng của hệ. Công thức nào sau đây là công thức tổng quát của nguyên lí thứ I của nhiệt động lực học?
\choice
{$A+\rm{Q}=\rm{0}{\rm}$}
{$\Delta U=\rm{Q}{\rm}$}
{$\Delta U=\rm{A}{\rm}$}
{\True $\Delta U=A+\rm{Q}{\rm}$}
\loigiai{}
\end{ex}
\begin{ex}
Phát biểu nào sau đây nói về nội năng là \textbf{không} đúng?
\choice
{Nội năng là một dạng năng lượng}
{\True Nội năng là nhiệt lượng}
{Nội năng của một vật có thể tăng lên hoặc giảm đi}
{Nội năng có thể chuyển hóa thành các dạng năng lượng khác}
\loigiai{}
\end{ex}
\begin{ex}
Khi nói về quá trình truyền nhiệt lượng khi cho hai vật tiếp xúc với nhau. Phát biểu nào sau đây là \textbf{sai}?
\choice
{\True Năng lượng nhiệt được truyền từ vật có nội năng lớn hơn sang vật có nội năng nhỏ hơn}
{Vật nóng hơn sẽ giảm nhiệt độ, vật lạnh hơn sẽ tăng nhiệt độ}
{Khi hai vật ở cùng nhiệt độ, không có truyền năng lượng nhiệt giữa chúng}
{Năng lượng nhiệt được truyền từ vật nóng hơn sang vật lạnh hơn}
\loigiai{}
\end{ex}
\begin{ex}
Khi các nguyên tử, phân tử cấu tạo nên vật chuyển động nhanh lên thì đại lượng nào sau đây tăng lên?
\choice
{Thế năng của vật tăng lên}
{Khối lượng của vật}
{Động năng của vật tăng lên}
{\True Nhiệt độ của vật}
\loigiai{}
\end{ex}
\begin{ex}
Khi nói về thang đo nhiệt độ Kelvin và Celsius, phát biểu nào sau đây là \textbf{sai}?
\choice
{Mối liên hệ về các giá trị nhiệt độ giữa hai thang đo là: $T(K) = t(^{\circ}C) + 273,15$}
{Nhiệt độ trong thang nhiệt độ Kelvin được kí hiệu là T, có đơn vị K}
{Nhiệt độ trong thang nhiệt độ Celsius được kí hiệu t, có đơn vị $^{\circ}C$}
{\True Một độ chia trên thang nhiệt độ Kelvin có giá trị gấp $273$ lần một độ chia trên thang nhiệt độ Celsius}
\loigiai{}
\end{ex}
\begin{ex}
Thể tích của một lượng khí khi bị nung nóng đã tăng thêm $0,02\ m^3$, còn nội năng tăng một lượng $4280\ J$. Xem như áp suất trong suốt quá trình trên là không đổi và bằng $1,5.10^5\ Pa$. Nhiệt lượng đã truyền cho khối khí trong quá trình trên là
\choice
{$1280\ J$}
{$4280\ J$}
{$3000\ J$}
{\True $7280\ J$}
\loigiai{}
\end{ex}
\begin{ex}
\immini[thm]{Đồ thị hình bên thể hiện quá trình tăng nhiệt độ theo thời gian của một chất rắn kết tinh khi được nung nóng. Nhiệt độ nóng chảy của chất rắn là
\choice[2]
{$210^\circ C$}
{$100^\circ C$}
{$320^\circ C$}
{\True $232^\circ C$}
}
{\begin{tikzpicture}[draw=Mapcolor, very thick, scale = 0.4, line join=round, line cap=round, >=stealth]
\draw[->, line width = 1.5] (0,0)--(8,0) node[above left] {\Large $\tau$};
\draw[->, line width = 1.5] (0,0)--(0,8) node[right] {\large t ($^\circ C$)};
\draw[fill = black] (0,0) circle(2pt) node [left] {\Large O};
\draw[-,dotted, line width = 1] (0,4.5) node [left]{$232$} -- (6,4.5);
\draw[line width = 2] (2,2)..controls+(75:3)and+(180:2.5)..(4.5,4.5)
..controls+(0:2.5)and+(-105:3)..(7,7);
\draw[-,dotted, line width = 1] (0,7) node [left]{$320$} -- (7,7);
\draw[-,dotted, line width = 1] (0,2) node [left]{$100$} -- (2,2);
\end{tikzpicture}
}
\loigiai{}
\end{ex}
\subsubsection*{PHẦN 2. Thí sinh trả lời từ câu 1 đến câu 4. Trong mỗi ý a), b), c), d) ở mỗi câu, thí sinh chọn đúng hoặc sai.}
\setcounter{ex}{0}
\begin{ex}%[2D1V1-1]
\immini[thm]{Khi nung nóng một khối khí chứa trong một xilanh có pit-tông đóng kín làm cho nhiệt độ của khối khí tăng. Pit-tông này có thể dịch chuyển không ma sát trong xilanh.}
{\begin{tikzpicture}[draw=Mapcolor, very thick, scale = 0.6, line join=round, line cap=round, >=stealth]% ĐÈN CỒN
\fill[blue!60] (-1,1) -- (1,1) -- (0.5,0) -- (-0.5,0) -- (-1,1);% CỒN
\draw[-,thick, line width = 2] (0,1.5) -- (-0.25,1.5) -- (-0.5,1) -- (-1,1) -- (-0.5,0) -- (0.5,0) -- (1,1) -- (0.5,1) -- (0.25,1.5) -- (0,1.5);% BÌNH ĐÈN CỒN
\draw[line width = 2.5] (0,1.65)..controls+(-90:0.5)and+(90:0.5)..(0,1.25)
..controls+(-35:1)and+(135:0.2)..(-0.35,0.5)..controls+(-35:0.2)and+(175:0.4)..(0.35,0.1);% DÂY
\fill[red] (0,1.65)..controls+(170:0.4)and+(-120:0.2)..(0,2.5)
..controls+(-100:0.2)and+(20:0.5)..(0,1.65);% LỬA
\draw[xshift = -1cm,yshift = 2.5cm,-,thick, line width = 6] (2.2,0.75) -- (3.5,0.75) -- (3.5,1) -- (3.5,0.5);
\fill[xshift = -1cm,yshift = 2.5cm,black!70](2,0) rectangle (2.4,1.5);
\draw[xshift = -1cm,yshift = 2.5cm,-,thick, line width = 3,rounded corners=2pt] (3,0) -- (0,0) -- (0,1.5) -- (3,1.5);
\fill[xshift = -1cm,yshift = 2.5cm,pattern=crosshatch dots,-,thick, line width = 2,rounded corners=4pt] (2,0) -- (0,0) -- (0,1.5) -- (2,1.5);
\end{tikzpicture}
}
\choiceTFt
{\True Nhiệt độ của khối khí trong xilanh thay đổi do quá trình truyền nhiệt}
{\True Khi cho pit-tông chuyển động tự do để khối khí dãn nở với áp suất không đổi, lúc này khối khí đã sinh công}
{\True Khi giữ pit-tông để thể tích khí không đổi thì toàn bộ nhiệt lượng khối khí nhận được dùng để tăng nội năng của khí (khối khí không tỏa nhiệt)}
{Khi nhận nhiệt, động năng của các phân tử là không đổi}
\loigiai{
}
\loigiai{}
\end{ex}
\begin{ex}%[2D1V1-1]
Giả sử một học sinh tạo ra một nhiệt kế sử dụng một thang nhiệt độ mới cho riêng mình, gọi là thang nhiệt độ $Z$, có đơn vị là ${ }^{\circ} \mathrm{Z}$. Trong đó, nhiệt độ của nước đá đang tan ở $1\ atm$ là $x^{\circ} \mathrm{Z}$ và nhiệt độ nước sôi ở $1\ atm$ là $y^{\circ} \mathrm{Z}$. Từ vạch $x^{\circ} \mathrm{Z}$ đến vạch $y^{\circ} \mathrm{Z}$ được chia thành $180$ khoảng, mỗi khoảng ứng với $1^{\circ} \mathrm{Z}$.
\choiceTFt
{Một độ chia trên thang nhiệt độ $Z$ bằng $1,8$ lần độ chia trên thang nhiệt độ Celsius}
{\True Mối liên hệ giữa x và y là: $y = x + 180$}
{Độ biến thiên nhiệt độ $18^\circ C$ trong thang nhiệt độ Celsius bằng với độ biến thiên nhiệt độ $10^\circ Z$ trong thang nhiệt độ Z}
{\True Nếu nhiệt độ cơ thể người là $37^\circ C$ tương ứng với $86,6^\circ Z$ thì giá trị của x là $20$}
\loigiai{
}
\loigiai{}
\end{ex}
\begin{ex}%[2D1V1-1]
\immini[thm]{Khi tiến hành nung nóng một chất rắn kết tinh bằng một bếp có công suất không đổi. Bỏ qua sự mất mát nhiệt lượng ra môi trường. Kể từ lúc bắt đầu đun người ta ghi nhận được đồ thị sự phụ thuộc của nhiệt độ của khối chất và thời gian đun như hình bên.}
{\begin{tikzpicture}[draw=Mapcolor, very thick, scale = 0.5, line join=round, line cap=round, >=stealth]
\draw[->, line width = 1.2] (0,0)--(8.5,0) node[above] {\large $\tau$ (phút)};
\draw[->, line width = 1.2] (0,0)--(0,4.5) node[right] {t($^o$C)};
\draw[fill = black] (0,0) circle(2pt) node [below left] {\large 0};
\draw[-,thick, line width = 1.5] (0,1) -- (3.57,3) -- (5,3) -- (7,4);
\draw[-,dashed, line width = 1] (0,3) node [left]{\large 327} -- (3.57,3) -- (3.57,0) node [below]{\large 39};
\draw[-,dashed, line width = 1] (5,3) -- (5,0) node [below]{\large 64};
\draw (0,1) node [left]{\large 27};
\draw[fill = black] (0,1) circle(2pt);
\draw[fill = black] (3.57,0) circle(2pt);
\draw[fill = black] (0,0.8) node [right] {\large A};
\draw[fill = black] (3.57,3) circle(2pt) node [above] {\large B};
\draw[fill = black] (5,3) circle(2pt) node [above] {\large C};
\draw[fill = black] (7,4) circle(2pt) node [above] {\large D};
\end{tikzpicture}
}
\choiceTFt
{\True Nhiệt độ nóng chảy của chất rắn trên là $327^\circ C$}
{\True Kể từ lúc bắt đầu đun, nhiệt lượng cần để chất rắn tăng lên đến nhiệt độ nóng chảy gấp $1,56$ lần nhiệt lượng cần cung cấp trong suốt giai đoạn nóng chảy}
{Đoạn AB trên đồ thị thể hiện quá trình chất rắn đang nóng chảy}
{Tại phút thứ $39$ chất rắn đã nóng chảy hoàn toàn}
\loigiai{
}
\loigiai{}
\end{ex}
\newpage
\begin{ex}%[2D1V1-1]
Nghiên cứu mô hình động học phân tử về cấu tạo chất.
\choiceTFt
{\True Mô hình động học phân tử được xây dựng dựa trên quan điểm là các chất có cấu tạo gián đoạn}
{Khi nhiệt độ cao các phân tử sẽ chuyển động, khi nhiệt độ thấp các phân tử sẽ đứng yên}
{\True Các chất được cấu tạo từ các hạt riêng biệt gọi là phân tử}
{\True Giữa các phân tử có lực hút và lực đẩy gọi chung là lực liên kết phân tử}
\loigiai{}
\end{ex}
\subsubsection*{PHẦN 3. Thí sinh trả lời từ câu 1 đến câu 6.}
\setcounter{ex}{0}
\begin{ex}
Một lượng khí được truyền $10\ kJ$ nhiệt năng để nóng lên, đồng thời khối khí giãn nở và thực hiện một công $8\ kJ$. Độ biến thiên nội năng của khối khí là bao nhiêu $kJ$?
\shortans[oly]{2}
\loigiai{}
\end{ex}

\begin{ex}
Một số nước trên thế giới sử dụng thang đo nhiệt độ Fahrenheit. Trong thang nhiệt này (ở áp suất tiêu chuẩn) nhiệt độ của nước đá đang tan là $32^\circ F$, của nước đang sôi là $212^\circ F$. Công thức chuyển đổi giữa thang đo Fahrenheit và thang đo Celsius là: $t(^\circ F) = 32 + 1,8.t(^\circ C)$. Nhiệt độ bằng bao nhiêu thì giá trị nhiệt độ trên hai thang đo là bằng nhau?
\shortans[oly]{-40}
\loigiai{}
\end{ex}

\begin{ex}
\immini[thm]{Hình bên là đồ thị biểu diễn sự phụ thuộc của nhiệt độ vào thời gian đun một ấm nước ở áp suất tiêu chuẩn. Nếu nhiệt lượng mà bếp tỏa ra không thay đổi trong suốt thời gian đun thì sau bao nhiêu giây kể từ lúc bắt đầu đun nước sẽ sôi?}
{\begin{tikzpicture}[draw=Mapcolor, very thick, scale = 0.75, line join=round, line cap=round, >=stealth]
\draw[gray!100,dashed] (0,0) grid[xstep=0.5,ystep=0.5] (3,3);
\draw[->, line width = 1] (0,0)--(3.5,0) node[above] {$\tau \,(s)$};
\draw[->, line width = 1] (0,0)--(0,3.5) node[right] {$t\,(^\circ )$};
\draw[fill = black] (0,0) circle(2pt) node [below left] {O};
\draw[-,thick, line width = 1] (0,1.5) -- (3,3);
\draw (0,1.5) node [left]{15};
\draw (2,0) node [below]{40};
\draw[fill = black] (0,1.5) circle(2pt);
\draw[fill = black] (2,2.5) circle(2pt);
\draw[fill = black] (2,0) circle(2pt);
\end{tikzpicture}
}
\shortans[oly]{340}
\loigiai{}
\end{ex}

\begin{ex}
Nhiệt lượng tỏa ra hay thu vào của một khối lượng chất $m$ trong quá trình truyền nhiệt được cho bởi hệ thức: $Q = mc\Delta t$; trong đó $c$ là một hằng số phụ thuộc vào cấu tạo của chất, $\Delta t$ là độ thay đổi nhiệt độ. Biết nhiệt độ ban đầu của nước sôi là $100^\circ C$, và của nước lạnh là $20^\circ C$. Bỏ qua sự trao đổi nhiệt với bình và môi trường.Nhiệt độ của hỗn hợp nước "3 sôi, 2 lạnh" sau khi có sự cân bằng nhiệt là bao nhiêu $^\circ C$?
\shortans[oly]{68}
\loigiai{}
\end{ex}

\begin{ex}
\impicinpar{Bằng các nghiên cứu, người ta phát hiện ra rằng các nguyên tử của nguyên tố X sắp xếp tuần hoàn tạo thành mạng tinh thể gồm các ô hình lập phương giống hệt nhau xếp chồng lên nhau (Hình a). Ở mỗi ô lập phương nhỏ nhất (gọi là ô mạng cơ sở) có một nguyên tử nằm tại tâm và ở mỗi đỉnh của nó đều có một nguyên tử (Hình b). Biết rằng chiều dài cạnh của mỗi ô lập phương cơ sở là $a = 2,87.10^{-10}\ m$. Biết khối lượng mỗi nguyên tử X là $9,3.10^{-26}\ kg$. Khối lượng riêng của nguyên tố X là bao nhiêu $kg/m^3$ (làm tròn kết quả đến hàng đơn vị)?}
{\begin{tikzpicture}[draw=Mapcolor, very thick, scale = 0.6, line join=round, line cap=round, >=stealth]
\draw[-,thick, line width = 1] (0,0) -- (4,0);
\draw[-,dashed, line width = 1] [xshift=0.75cm, yshift=0.5cm] (0,0) -- (4,0);
\draw[-,dashed, line width = 1] [xshift=1.5cm, yshift=1cm] (0,0) -- (4,0);
\draw[-,thick, line width = 1] (0,2) -- (4,2);
\draw[-,dashed, line width = 1] [xshift=0.75cm, yshift=.5cm] (0,2)  -- (4,2);
\draw[-,dashed, line width = 1] [xshift=1.5cm, yshift=1cm] (0,2) -- (4,2);
\draw[-,thick, line width = 1] [] (0,4) -- (4,4);
\draw[-,thick, line width = 1] [xshift=0.75cm, yshift=0.5cm] (0,4) -- (4,4);
\draw[-,thick, line width = 1] [xshift=1.5cm, yshift=1cm] (0,4) -- (4,4);
\draw[-,thick, line width = 1] (0,0) -- (0,4);
\draw[-,dashed, line width = 1] [xshift=0.75cm, yshift=0.5cm] (0,0) -- (0,4);
\draw[-,dashed, line width = 1] [xshift=1.5cm, yshift=1cm] (0,0) -- (0,4);
\draw[-,thick, line width = 1] (2,0) -- (2,4);
\draw[-,dashed, line width = 1] [xshift=0.75cm, yshift=0.5cm] (2,0) -- (2,4);
\draw[-,dashed, line width = 1] [xshift=1.5cm, yshift=1cm] (2,0) -- (2,4);
\draw[-,thick, line width = 1] (4,0) -- (4,4);
\draw[-,thick, line width = 1] [xshift=0.75cm, yshift=0.5cm] (4,0) -- (4,4);
\draw[-,thick, line width = 1] [xshift=1.5cm, yshift=1cm] (4,0) -- (4,4);
\draw[-,dashed, line width = 1] (0,0) -- (1.5,1);
\draw[-,dashed, line width = 1] [xshift=0cm, yshift=2cm] (0,0) -- (1.5,1);
\draw[-,thick, line width = 1] [xshift=0cm, yshift=4cm] (0,0) -- (1.5,1);
\draw[-,dashed, line width = 1] (2,0) -- (3.5,1);
\draw[-,dashed, line width = 1] [xshift=0cm, yshift=2cm] (2,0) -- (3.5,1);
\draw[-,thick, line width = 1] [xshift=0cm, yshift=4cm] (2,0) -- (3.5,1);
\draw[-,thick, line width = 1] (4,0) -- (5.5,1);
\draw[-,thick, line width = 1] [xshift=0cm, yshift=2cm] (4,0) -- (5.5,1);
\draw[-,thick, line width = 1] [xshift=0cm, yshift=4cm] (4,0) -- (5.5,1);
\draw[line width = 1, fill = black] (0,0) circle(5pt);
\draw[line width = 1, fill = black] (2,0) circle(5pt);
\draw[line width = 1, fill = black] (4,0) circle(5pt);
\draw[line width = 1, fill = black] (0,2) circle(5pt);
\draw[line width = 1, fill = black] (2,2) circle(5pt);
\draw[line width = 1, fill = black] (4,2) circle(5pt);
\draw[line width = 1, fill = black] (0,4) circle(5pt);
\draw[line width = 1, fill = black] (2,4) circle(5pt);
\draw[line width = 1, fill = black] (4,4) circle(5pt);
\draw[line width = 1, fill = black] (0.75,0.5) circle(5pt);
\draw[line width = 1, fill = black] (2.75,0.5) circle(5pt);
\draw[line width = 1, fill = black] (4.75,0.5) circle(5pt);
\draw[line width = 1, fill = black] (0.75,2.5) circle(5pt);
\draw[line width = 1, fill = black] (2.75,2.5) circle(5pt);
\draw[line width = 1, fill = black] (4.75,2.5) circle(5pt);
\draw[line width = 1, fill = black] (0.75,4.5) circle(5pt);
\draw[line width = 1, fill = black] (2.75,4.5) circle(5pt);
\draw[line width = 1, fill = black] (4.75,4.5) circle(5pt);
\draw[line width = 1, fill = black] (1.5,1) circle(5pt);
\draw[line width = 1, fill = black] (3.5,1) circle(5pt);
\draw[line width = 1, fill = black] (5.5,1) circle(5pt);
\draw[line width = 1, fill = black] (1.5,3) circle(5pt);
\draw[line width = 1, fill = black] (3.5,3) circle(5pt);
\draw[line width = 1, fill = black] (5.5,3) circle(5pt);
\draw[line width = 1, fill = black] (1.5,5) circle(5pt);
\draw[line width = 1, fill = black] (3.5,5) circle(5pt);
\draw[line width = 1, fill = black] (5.5,5) circle(5pt);
\draw (2,-1) node []{Hình a};
%%%%%%%%%%%%%%%%%%
\draw[xshift = 8 cm, line width = 1, fill = black] (0,0) circle(5pt);
\draw[xshift = 8 cm,line width = 1, fill = black] (3,0) circle(5pt);
\draw[xshift = 8 cm,line width = 1, fill = black] (0,3) circle(5pt);
\draw[xshift = 8 cm,line width = 1, fill = black] (3,3) circle(5pt);
\draw[xshift = 8 cm,line width = 1, fill = black] (1,1.3) circle(5pt);
\draw[xshift = 8 cm,line width = 1, fill = black] (4,1.3) circle(5pt);
\draw[xshift = 8 cm,line width = 1, fill = black] (1,4.3) circle(5pt);
\draw[xshift = 8 cm,line width = 1, fill = black] (4,4.3) circle(5pt);
\draw[xshift = 8 cm,-,thick, line width = 1] (0,0) -- (0,3) -- (3,3) -- (3,0) -- (0,0);
\draw[xshift = 8 cm,-,thick, dashed = 1] (0,0) -- (1,1.3) -- (1,4.3);
\draw[xshift = 8 cm,-,thick, dashed = 1] (1,1.3) -- (4,1.3);
\draw[xshift = 8 cm,-,thick, line width = 1] (0,3) -- (1,4.3) -- (4,4.3) -- (4,1.3) -- (3,0);
\draw[xshift = 8 cm,-,thick, line width = 1] (3,3) -- (4,4.3);
\draw[xshift = 8 cm,line width = 1, fill = black] (2,2.15) circle(5pt);
\draw[xshift = 8 cm,-,thick, dashed = 1] (1,4.3) -- (3,0);
\draw[xshift = 8 cm,-,thick, dashed = 1] (0,3) -- (4,1.3);
%\draw[-,thick, dashed = 1] (0,0) -- (4,4.3);
%\draw[-,thick, dashed = 1] (1,1.3) -- (3,3);
\draw (9.5,-1) node []{Hình b};
\draw (10.5,4.75) node []{a};
\draw (12.5,3) node []{a};
\draw (5.5,-2) node []{\textit{Trên hình a không thể hiện các }};
\draw (5.5,-3) node []{\textit{nguyên tử ở tâm của các ô mạng}};
\end{tikzpicture}
}
\shortans[oly]{7868}
\loigiai{
\begin{itemize}
\item  Xét ô lập phương cơ sở như hình b
\item Mỗi nguyên tử tại mỗi đỉnh đóng góp $\dfrac{1}{8}$ khối lượng của mình cho khối lượng của ô tinh thể.
\item Nguyên tử ở tâm đóng góp toàn bộ khối lượng của mình cho khối lượng ô tinh thể.
\item Vậy khối lượng ô lập phương cơ sở là: $M = 8.\dfrac{1}{8}.m + m  = 2m$.
\item Thể tích mỗi ô lập phương: $V=a^3$.
\item Khối lượng riêng của chất X: $D = \dfrac{M}{V} = \dfrac{2m}{a^3} = 7868\ kg/m^3$.
\end{itemize}
}
\loigiai{}
\end{ex}

\begin{ex}
Một nhiệt kế gồm phần vỏ thủy tinh và phần chất lỏng bằng rượu. Biết rằng khi nhiệt kế chỉ $20^\circ C$ thì phần chất lỏng trong nhiệt kế có thể tích là $V_0$. Khi nhiệt kế chỉ $40^\circ C$ thì phần rượu trong nhiệt kế có thể tích $V_1=1,12V_0$. Khi nhiệt kế chỉ $80^\circ C$ thì phần rượu trong nhiệt kế có thể tích $V_2$. Bỏ qua sự nở vì nhiệt của phần vỏ thủy tinh. Tỉ số $\dfrac{V_2}{V_0}$ có giá trị bằng bao nhiêu (làm tròn kết quả đến hàng phần trăm)?
\shortans[oly]{1,36}
\loigiai{}
\end{ex}




